\section{Follow-up implementation informed by the literature}
\label{sec:followup}
The literature review and the first experiment changed the optimization target.
The earlier selector saved inner-evaluator work while transaction preparation
dominated several complete updates. Changing-rule maintenance already has a
substantial history~\cite{volz2003maintenance,kotowski2011,xu2023}; recent join
research emphasizes the costs of execution choices and filtering
\cite{qiao2026}. We therefore implemented three further specializations and
evaluated them against both preserved commits and a same-source control.
These are engineering adaptations of established ideas. They neither implement
RPT+ nor replace B/F$^c$, modular maintenance, or a worst-case-optimal join.

\subsection{Reusing validation for monotone transactions}
The engine records the validated rule sequence, assertion set, and predicate
arities. An insertion-only transaction may reuse this certificate only when
the public rule and assertion containers still match it exactly. New facts and
rules are checked against the recorded arities and one another before the
transaction changes state. Submitted removals, a missing certificate, or direct
container mutation trigger complete validation. Deletion can release the last
occurrence fixing a predicate's arity, so applying the monotone shortcut there
would be unsound.

\begin{proposition}
For an unchanged validated input state whose embedded terms and identifiers
have stable values, equality, and hashes, checking the added facts and rules
against its arity map preserves the engine's input-validation conditions for
insertion-only transactions.
\end{proposition}
\begin{proof}
Groundness and permitted-term conditions on old facts, and rule safety and
permitted-term conditions on old rules, continue to hold because those objects
retain their structure and embedded values. These local conditions are checked on
every genuinely new object. The remaining shared condition is that every
occurrence of a predicate has the same arity. Extending a copy of the old arity
map while checking all new occurrences enforces that condition across old and
new objects and within the additions. The extended certificate is installed
only after these checks succeed.
\end{proof}

The API's frozen rule and atom records are shallow: arbitrary custom hashable
values can still have mutable internals. The stated stable-value assumption is
required for both the certificate and ordinary hash indexes. RDF terms and the
primitive identifiers used in these experiments satisfy it; snapshot equality
does not detect arbitrary in-place changes inside custom value objects.

When no rule is genuinely new, prepared plans are retained. A complete old
materialization has already registered its domain, including the private
nonempty-domain seed. Insertion of new rules therefore requires explicit
domain seeding only for their previously unseen constants; ingesting new facts
registers their own terms. Incomplete old states retain rematerialization.
Snapshot equality, assertion-set construction, and immutable snapshots still
incur work proportional to the existing state. This is a reduction in repeated
Python validation and preparation, not an $O(|\Delta|)$ bound for the complete
transaction.

\subsection{Maintaining surviving DRed indexes}
After overdeletion has finished using the old closure, ordinary DRed rebuilds
the index over the survivors. The candidate can instead erase each overdeleted
tuple from its relation and every corresponding column bucket. It uses erasure
when $2|D|\leq|I\setminus D|$, where $D$ is the overdeleted set and $I$ the old
indexed closure, and otherwise rebuilds. Empty buckets and relations are
removed. The threshold is an implementation heuristic, not a competitive
runtime guarantee.

For each relation $P$, column $j$, and value $v$, the index invariant is
\[
 B_{P,j,v}=\{t\in I_P:t_j=v\}.
\]
Erasing every $d\in D_P$ from $I_P$ and from $B_{P,j,d_j}$ leaves exactly the
buckets obtained by rebuilding over $I_P\setminus D_P$. Deletion occurs only
after the overdelete traversal, so its old-index access remains valid.
Rederivation then sees the same survivor index under either implementation.
The eligibility checks for equality, function terms, and complete old closures
are unchanged. This optimization reduces index construction costs; it does not
reduce the logical overdeletion set or solve incremental equality splitting.

\subsection{Positional bindings for ordinary positive joins}
For positive bodies without equality or inequality atoms, a reusable plan
assigns each variable a position in a binding vector. Atom terms are compiled
to variable positions or constant positions. Each recursive step retains the
existing dynamic choice of the smallest available exact index bucket, extends
a copied binding vector, and checks repeated variables. Constants are
canonicalized on every invocation because equality representatives can change.
The final vector is converted back to the original variable-binding interface.
Unary same-variable plans and the generic fallback remain available.

For validated bodies, the preservation argument is an induction on the remaining
body atoms; nonground Skolem patterns in bodies are rejected by validation. A
vector position represents exactly the binding of its corresponding variable.
The chosen index bucket is a superset of matches for the current partial
binding; constant and repeated-variable checks retain precisely the compatible
tuples. Every compatible extension recurses on the remaining atoms. At an
empty body, conversion preserves all bindings. Semi-naive evaluation keeps the
designated delta \emph{occurrence}, including when several atoms share a
predicate. Changing atom order does not change the set of substitutions.
This specializes an interpreter; it adds neither semijoin filtering nor the
Bloom-filter pipelines of Qiao et al.~\cite{qiao2026}.

Before the external timing series, the complete suite passed 698 tests. New
controls include 420 randomized insertion states, 140 randomized mixed
transactions with index comparisons, and 180 generated relational bodies
compared with an independent Cartesian-product evaluator. Tests also cover
failed-validation atomicity, public-container mutation, new domain constants,
equality fallback, and repeated delta predicates. These checks address the
implemented contracts rather than full OWL conformance.

\subsection{Matched author-program and SQL-component protocols}
The four local configurations are immutable commits \code{b1254c4} and
\code{5531be4}, the new source with all three switches disabled, and that same
source with all three enabled. The earlier adaptive unary selector and support
certificates remain disabled throughout. The same-source control distinguishes
the enabled bundle from changes that affect both paths. Comparing the bundle
does not isolate the contribution of each maintenance optimization.

The ZodiacEdge author repository~\cite{zodiacartifact} supplies a 128-rule
positive LUBM program and two marked rule subsets: eight rules marked $*$ and
sixteen marked $\#$, with one rule shared. Each scenario materializes without
the marked rules, inserts them, and removes them. Official UBA generation
provides 100,543 ABox facts after excluding ontology metadata. The author's
generation seed and file hashes are unavailable, so matching its published
input count is not proof of byte-identical data. The final author-program
closure has 238,474 non-domain tuples, including the 100,543 wrapped input
facts. It is a different semantic artifact from full-ontology DLP L3 reasoning.

Four balanced process blocks evaluate all four configurations. Ten fixed,
seeded samples each remove and reinsert 1,000 assertions under the complete
program. This follows the small-update scale used by Hu, Motik, and
Horrocks~\cite{hu2018}, but does not reproduce their Reactome or UOBM programs.
Thirteen fresh naive reference closures cover the full program, both partial
programs, and the ten fact-deletion states. Every timed worker checks complete
typed-tuple digests after every transaction outside the timer. The ten samples
within one worker are correlated: forty paired sample ratios are not forty
independent process replicates. Reported summaries are descriptive.

The Qiao-related component experiment uses official DuckDB \code{dbgen} data
at scale factors 0.01 and 0.1 and the positive joins underlying TPC-H Q3, Q5,
and Q9. It retains the original tables and filters, but projects every
participating row's complete primary key. Under the checked key constraints,
that projection identifies the entire tuple of input rows injectively and
therefore preserves the pre-aggregation SQL bag's multiplicity. Exact answers
are checked against independently executed SQL over the original database.
The engine's rule-free base closure is independently checked as its input
facts plus all required domain facts.

DuckDB performs date/string selections and projection; their time and complete
transfer to Python are recorded. Fact conversion, engine construction/indexing,
and complete join consumption are separate intervals, all included in the cold
component total. The three components run sequentially inside each fresh
worker, with a new reasoner per component but shared DuckDB and filesystem
caches. Aggregation, arithmetic, ordering, and top-$k$ are outside
this experiment. Three blocks, two scales, and four configurations yield 24
fresh workers and 72 component observations. Three rotations do not completely
balance four configurations. These small components are useful regression and
optimization probes, not a full TPC-H result or a reproduction of Qiao et al.'s
SF100 evaluation.

\subsection{Follow-up measurements}
\begin{table}[ht]
\caption{Author-supplied ZodiacEdge LUBM program: four fresh process blocks per configuration. Median seconds; initial includes engine construction and materialization. Parsing and outcome validation are outside these intervals.}
\label{tab:zodiac}
\centering\small
\begin{tabular}{@{}llrrr@{}}
\toprule
Subset & Configuration & Initial & Insert rules & Delete rules \\
\midrule
$\ast$ & b1254c4 & 3.5382 & 0.3338 & 2.3322 \\
$\ast$ & 5531be4 & 3.2811 & 0.3343 & 2.1418 \\
$\ast$ & current-fixed & 2.9002 & 0.3182 & 1.9661 \\
$\ast$ & candidate & 2.7809 & 0.1020 & 1.3576 \\
$\#$ & b1254c4 & 3.2615 & 0.4923 & 1.5962 \\
$\#$ & 5531be4 & 3.1646 & 0.4933 & 1.5541 \\
$\#$ & current-fixed & 2.7764 & 0.4403 & 1.4061 \\
$\#$ & candidate & 2.4551 & 0.2458 & 1.0838 \\
\bottomrule
\end{tabular}
\end{table}

\begin{table}[ht]
\caption{Supplemental 1,000-fact updates on the same author program. Median seconds across ten fixed subsets in each of four process blocks. These are forty correlated case/block observations per configuration.}
\label{tab:zodiac-facts}
\centering\small
\begin{tabular}{@{}lrr@{}}
\toprule
Configuration & Delete facts & Reinsert facts \\
\midrule
b1254c4 & 2.3727 & 0.1349 \\
5531be4 & 2.1627 & 0.1210 \\
current-fixed & 1.8719 & 0.1141 \\
candidate & 1.4258 & 0.0395 \\
\bottomrule
\end{tabular}
\end{table}

\begin{table}[ht]
\caption{Fully consumed join kernel for TPC-H-derived Q3/Q5/Q9 components (ms). Three fresh workers per scale/configuration. Final column is the median paired candidate/current-fixed ratio; smaller favors the candidate.}
\label{tab:tpch-joins}
\centering\small
\begin{tabular}{@{}llrrrrr@{}}
\toprule
SF & Join & b1254c4 & 5531be4 & Fixed & Candidate & Ratio \\
\midrule
0.01 & Q3 & 5.03 & 4.76 & 4.69 & 3.69 & 0.755 \\
0.01 & Q5 & 20.26 & 21.30 & 21.26 & 11.34 & 0.538 \\
0.01 & Q9 & 35.98 & 34.76 & 34.43 & 23.08 & 0.670 \\
0.1 & Q3 & 45.10 & 44.80 & 43.05 & 31.78 & 0.752 \\
0.1 & Q5 & 941.87 & 957.38 & 949.31 & 500.09 & 0.523 \\
0.1 & Q9 & 390.87 & 378.17 & 380.53 & 266.40 & 0.700 \\
\bottomrule
\end{tabular}
\end{table}

\begin{table}[ht]
\caption{Cold component total for TPC-H-derived Q3/Q5/Q9 components (ms). Three fresh workers per scale/configuration. Final column is the median paired candidate/current-fixed ratio; smaller favors the candidate.}
\label{tab:tpch-total}
\centering\small
\begin{tabular}{@{}llrrrrr@{}}
\toprule
SF & Join & b1254c4 & 5531be4 & Fixed & Candidate & Ratio \\
\midrule
0.01 & Q3 & 224.58 & 221.31 & 226.80 & 223.30 & 0.985 \\
0.01 & Q5 & 428.23 & 429.61 & 438.62 & 419.32 & 0.974 \\
0.01 & Q9 & 657.89 & 650.38 & 654.30 & 627.63 & 0.959 \\
0.1 & Q3 & 4,040.74 & 4,130.65 & 3,999.45 & 4,057.00 & 1.008 \\
0.1 & Q5 & 8,767.87 & 8,767.37 & 8,712.05 & 8,112.88 & 0.919 \\
0.1 & Q9 & 11,740.80 & 11,615.00 & 11,533.75 & 10,994.71 & 0.952 \\
\bottomrule
\end{tabular}
\end{table}


\subsection{An actual external engine on the same host}
The comparator uses the pinned, unmodified ZodiacEdge Program API and parser.
Its unary atoms use physical RDF type triples, whereas our store uses unary
relations. An injective encoding retains RDF term identity, and both outputs
are decoded to the same typed tuples for complete closure checks. Native input
index creation is included in initial compute. Conversion is recorded separately
and excluded from both implementations' compute intervals. Native whole-process
RSS includes adapter allocations, so it is not an isolated engine-memory ranking.
The program has no predicate variables; fixed predicate labels and injectively
encoded data constants occupy separate namespaces. Translating unary atoms to
type triples and binary atoms to property triples therefore gives a bijection
between rule substitutions. This validates the adapter's semantic contract;
full-run checks still test the engines' actual computation.
\begin{table}[ht]
\caption{Matched comparison with the unmodified author engine: four alternating paired blocks, seconds. Initial includes native input indexing. Final column is the median paired candidate/native ratio; below one favors our candidate.}
\label{tab:zodiac-native}
\centering\small
\begin{tabular}{@{}llrrr@{}}
\toprule
Subset & Operation & ZodiacEdge & Candidate & Ratio \\
\midrule
$\ast$ & Initial & 1.4001 & 2.4567 & 1.756 \\
$\ast$ & Insert rules & 0.5590 & 0.0880 & 0.161 \\
$\ast$ & Delete rules & 0.7850 & 1.1965 & 1.522 \\
$\#$ & Initial & 1.6595 & 2.1807 & 1.318 \\
$\#$ & Insert rules & 0.0589 & 0.2138 & 3.653 \\
$\#$ & Delete rules & 0.0234 & 0.9709 & 41.335 \\
\bottomrule
\end{tabular}
\end{table}


\subsection{Interpreting the maintenance improvement}
All 16 ZodiacEdge-program workers pass all 416 recorded complete-closure stage
checks. Relative to \code{b1254c4}, median paired insertion ratios are 0.3144
for $*$ and 0.4807 for $\#$, corresponding to 3.18-fold and 2.08-fold speedups.
Deletion ratios are 0.6119 and 0.6969. These are new matched local comparisons
on the author's program, not ratios against its published README times.

The same-source control is especially useful for attribution. Candidate/fixed
rule-insertion ratios are 0.3265 and 0.5470; deletion ratios are 0.7240 and
0.8230. Every individual paired update ratio favors the candidate in both
scenarios. For initial compute, however, the $*$ median ratio is 0.9770 and
two of the four individual ratios exceed one. We do not interpret that small
aggregate difference as a reliable initialization benefit. The $\#$ initial
ratio is 0.9051, with all four observations below one. No predeclared statistical
significance test or broad default-adoption criterion was attached to this
follow-up study.

Across the ten fixed 1,000-fact subsets and four blocks, pooled matched
candidate/\code{b1254c4} ratios are 0.6251 for deletion and 0.2808 for
reinsertion. Against the same-source fixed control they are 0.7169 and 0.3064.
These medians describe the sampled workload, with cases nested within processes.
They do not supply forty independent repetitions or establish the performance
of Oxford's larger UOBM or Reactome workloads.

The candidate preserves the original logical DRed work contract. For example,
in the first block it erases 24,933 indexed tuples for $*$ rule deletion and
24,045 for $\#$, rather than rebuilding all survivor buckets. Insertion can
reuse validated old inputs and plans. The experiment tests the combined changes;
it does not independently assign an exact portion of each speedup to validation,
index erasure, or positional joins. The original adaptive unary strategy and
proof certificates are disabled, so these results do not reverse their earlier
default-adoption findings.

\subsection{Positive-join gains and remaining intermediate work}
All 72 TPC-H component observations pass exact SQL row-identity and complete
base-closure checks. Paired candidate/fixed join ratios range from 0.523 to
0.755 over the six scale/component cases. The corresponding component totals
range from 0.919 to 1.008. Large Q5 improves about 47.7\% in the join interval,
but only 8.1\% in the sum of extraction, encoding, indexing, and joining.
Large Q3 does not show a component-total gain. These differences matter because
the Python loading and representation costs dominate this adapter.

Candidate and fixed configurations enumerate identical numbers of candidate
rows in every case. Q5 increases from 13,739 to 621,901 candidates between
SF0.01 and SF0.1, a 45.27-fold expansion for approximately tenfold input growth.
Positional bindings reduce interpretation cost but leave this join order
unchanged. This observation motivates the separately versioned lookahead
experiment below; it does not by itself identify the causal ordering decision.

The original component worker executes three queries sequentially. Later-query
fact encoding was faster in the candidate despite identical encoding logic.
Inspection found recursive local generator closures retaining captured state
until cyclic garbage collection. The earlier summed phase intervals exclude
explicit teardown and validation, so they are not all application CPU costs.
This is a plausible interaction, not a proven causal explanation of those
timings. The follow-up breaks the closure cycle in both new arms and gives each
component its own process, while retaining the original observations unchanged.

\subsection{The external comparison is mixed}
All eight native-comparison workers pass their six full-closure checks, giving
48 checked states. Relative to the unmodified ZodiacEdge engine, our candidate's
paired $*$ insertion ratio is 0.161, approximately a 6.2-fold improvement.
For $*$ deletion, however, the ratio is 1.522; for $\#$ insertion it is 3.653,
and for $\#$ deletion it is 41.335. Initial computation also favors ZodiacEdge.
One native $\#$ insertion is an outlier at approximately 0.215 seconds, compared
with approximately 0.058 seconds for the other three; all observations remain
in the artifact. Four paired blocks justify descriptive comparisons on these
operations, not a broad statistical or state-of-the-art ranking.

The large $\#$ deletion gap warrants diagnosis before another algorithmic claim.
Index erasure reduces physical rebuilding but leaves transaction-wide
preparation passes in place. A separately instrumented run, described below,
attributes most of this case's cost to those passes, while its overdeletion
set is exactly the set of lost facts. Thus this particular gap should not be
attributed to unnecessary logical overdeletion. Dependency-local maintenance
and safe reuse of validation evidence are relevant next directions, informed
by ZodiacEdge~\cite{xu2023} and modular materialization~\cite{hu2022modular}.

\subsection{Diagnostic profile of the rule-deletion gap}
\label{sec:deletion-diagnostic}
After the paired native comparison, we profiled one deletion of the sixteen
$\#$-marked rules using CPython 3.12.9 \code{cProfile}. The process used the
frozen primary candidate, its three enabled optimizations, the same 100,543
assertions and 128-rule program, and hash seed 7300. It first materialized the
partial program and inserted the marked rules; only their subsequent removal
was profiled. All three full typed-tuple closure checks passed. This is a
separate instrumented diagnostic, not an additional timing replicate.

The profiled public update accumulated 2.130745 seconds. Four disjoint call
groups accounted for 1.553658 seconds, or 72.9\%: canonicalizing all 100,543
assertions for protection, 0.551646 seconds; two checks of old and new DRed
eligibility, 0.520192 seconds; full input validation, 0.295590 seconds; and
recomputing active-domain constants, 0.186231 seconds. These cumulative times
include descendants but not one another. Nested normalization and literal-key
operations, or the enclosing DRed routine, must not be added again. All other
update work accounts for the remaining 0.577087 seconds. The engine's internal
update counter excludes validation and eligibility checks and therefore covers
a narrower interval.

Only sixteen rules were evaluated. The deletion overdeleted 24,045 facts and
rederived none, exactly matching the difference between the complete closures
of 238,474 and 214,429 tuples. Thus this observation contains no excess logical
overdeletion. Global preparation and protection passes dominate the observed
candidate cost. Since the native engine was not profiled, this result does not
quantitatively decompose the separately measured 41.3-fold runtime ratio or
show that eliminating those passes would close the gap.

Maintaining auxiliary state incrementally is established practice in view and
materialization maintenance~\cite{gupta1993,motik2019}; dependency analysis also
limits reconsideration under rule changes~\cite{kotowski2011,xu2023}. Applying
these general principles to the engine's metadata suggests eligibility
certificates, predicate occurrence counts, and active-domain reference counts.
The exact proposals here are not claims about the implementations in those
papers. Any certificate must detect relevant public-state changes and handle
datatype aliases and equality history. Arity counts must release a predicate's
constraint when its last occurrence disappears; domain counts must handle
constants supported only by removed rules and preserve the nonempty-domain
seed. Within an eligible DRed update, identity representatives also suggest
avoiding repeated canonicalization of unchanged assertions while preserving
inequality symmetry and domain protection.

These changes remain proposals, not implemented or measured improvements.
Deterministic profiling perturbs runtime, particularly for call-heavy Python
paths, and this single transaction cannot establish general workload behavior
or an achievable speedup. The compact artifact
\code{benchmarks/zodiac-deletion-profile.json} records frozen source and
profiler source hashes, the raw profile hash, independent closure checks,
nonoverlapping costs, and the instrumented engine counters.


